\section{Compiler extensions and the normalized 246-state control}
\label{sec:compiler-extensions}

The source commutation proved earlier creates a different optimization
starting point.  To explore it, the compiler search was extended in three
ways: physical state assignments are transported through source rotations;
arithmetic terminals can reuse existing selector entries; and relative
control operands can be recomputed after a layout change.  These changes
enlarge the search space.  They do not by themselves constitute a proof
that an emitted binary table implements its source.

\subsection{Transporting physical assignments}

Write a shared source node as $N(A,B)$, and call the physical decoder state
assigned to that node its \emph{color}.  Different source nodes may have the
same color because the binary executor also uses information on the tape.
In addition, a row of a decoder color may execute arithmetic or control
work after its first decoding use.  Consequently, counting distinct source
nodes does not count physical states.

For each reassociation
\[
 N(A,N(B,C))\longleftrightarrow N(N(A,B),C),
\]
the search considers both ways of transporting the two local colors.
Ancestor signatures are recomputed by structural interning.  When two
signatures become equal, their colors may be merged if the protected-row
constraints permit it; an unused unprotected color can then be removed
from the emitted table.  All supporting states remain explicit, and every
nonhalting transition target must be defined.  A merger of protected
colors is rejected instead of being counted as a state saving.

The protection data include rows observed outside their first source
decoder use, targets reached from supporting states, and shared decoder
entry contexts.  These finite observations are search filters, not a
universal theorem about all executions.  The final proof obligation is
the separate symbolic execution of every source instruction frame.  In
particular, a candidate is not accepted merely because its structural
colors fit or its finite sample executions succeed.

For ordinary reassociations, a complete finite source check compares all
6,460 instruction occurrences and their guarded successor occurrences.
There are 8,390 successor edges and one halting case.  No allegedly
unreachable guard is omitted in this compiler lane.  The independently
checked color-106 commutation changes the preceding source program; later
reassociations are compared to that commuted source and compose with its
finite-region proof.

\subsection{Normalizing arithmetic terminals without adding states}

The earlier emitter retained a terminal's original borrowed handler where
possible.  That unnecessarily restricted some instruction addresses: a
borrowed handler may inspect preceding address bits even though a direct
selector entry already exists elsewhere in the same table.

The new emitter proposes the following existing entries.  Register numbers
in the table are logical source numbers, starting at zero; state labels
are the original physical labels, not the canonical renumbering.

\begin{center}
\begin{tabular}{@{}lrrrrrr@{}}
\toprule
Instruction & \multicolumn{6}{c}{Logical register $\longmapsto$ entry state}\\
\midrule
Increment & $1\mapsto70$ & $3\mapsto201$ & $4\mapsto169$
 & $5\mapsto72$ & $6\mapsto39$ & $7\mapsto93$\\
Decrement & $1\mapsto49$ & $3\mapsto21$ & $4\mapsto173$
 & $5\mapsto53$ & $6\mapsto25$ & \\
\bottomrule
\end{tabular}
\end{center}

At an unprotected terminal edge reading $b$, the proposed row is
$(b,R,q)$ for the listed entry $q$.  This uses an already counted state.
For these proposals the search removes the artificial requirement that
the last two address bits equal those of the corresponding original
occurrence.  Other arithmetic terminals remain constrained: some use a
common entry for increment and decrement, or borrow a decoder row whose
meaning depends on its entry context.  Existing protected rows are never
silently normalized.

The table of entries is an emission strategy, not a blanket theorem for
every conceivable prefix.  Its retained result is justified by the full
candidate-specific instruction certificates.  This separation matters:
relaxing a source suffix constraint can expose a valid layout, but can also
expose an invalid borrowed-row context.  Only the universal frame audit
distinguishes the two.

The principal retained result is
\texttt{machines/normalized246\_old.tm}, canonically renumbered as
\texttt{machines/normalized246.tm}.  Its source graph has 203 shared pairs,
6,460 leaves and maximum path depth 40.  It still has 246 working states.
Its literal-table SHA-256, before canonical renumbering, is
\begin{center}\small
\texttt{b2a2e6608a603e3b8dda9fc3c4195e6376}\\
\texttt{a934bba90519fa4d5e30306ed6ea25}.
\end{center}
The split above is typographical; the two lines concatenate to one hash.

Independent source verification proves that its instruction operands and
all resolved successors agree with the commuted seed by occurrence.
Independent blank-tape simulation reaches the full reference frame after
410,735 moves.  The two edited color-106 rows are used during this prefix,
so the proof compares the resulting complete frame rather than asserting
that the old transition trace is unchanged.

\subsection{A larger invariant gap costs no machine states}

The initial standard frame has head-to-header distance $d=181$.  Earlier
experiments frequently required an instruction theorem at $d=40$.  That
requirement can become unnecessarily restrictive after reassociation,
because an address may approach or exceed that depth.

For the retained control, the symbolic frame theorem is instead stated
for $d\geq64$.  The executor and its arithmetic rules are unchanged.  In
the inherited affine checker, the single expression specifying the
initial symbolic gap changes from $40-|p|$ to $64-|p|$; the requested final
gap is changed consistently.  The arbitrary additional nonnegative gap,
all natural register values and the arbitrary nonempty spare suffix remain
universally quantified.

\begin{lemma}[Preservation of the stronger frame domain]
Suppose the blank prefix reaches a standard frame with $d\geq64$ and the
source's maximum address depth is at most 64.  If each instruction contract
preserves the register encoding and changes the distance by a nonnegative
amount, every later source instruction begins inside the same frame domain.
\end{lemma}
\begin{proof}
The first instruction satisfies the domain by hypothesis.  The register
contracts preserve the standard finite-word form.  The distance increases
by one for a positive decrement and by one for an address overflow, when
the corresponding event occurs, and otherwise does not decrease.  Thus
$d'\geq d\geq64$ after each instruction.  Induction on the number of source
instructions proves the assertion.
\end{proof}

For the retained table, all 8,391 guarded symbolic cases pass at this
stronger minimum distance, together with the arbitrary-register arithmetic
and finite-list induction obligations.  Every certified instruction uses
at least 13 literal moves, which excludes an infinite source computation
being represented by only finitely many binary moves.  The nine-zero
scanner remains in the retained machine.

\subsection{What the scanner progress does and does not show}

The following source layouts all retain 246 binary states.  A decrement
whose address ends in seven zeros contributes two guarded scanner
obstructions; an increment can contribute one.  These counts concern the
specific attempted deletion of scanner state 74, not arbitrary possible
state reductions.

\begin{center}
\begin{tabular}{@{}lrrrr@{}}
\toprule
Layout & Source pairs & Maximum depth & Long-tail decrements
 & Guarded obstructions\\
\midrule
\texttt{s3\_best\_010} & 202 & 39 & 12 & 25\\
\texttt{s7\_best\_018} & 203 & 43 & 10 & 21\\
\texttt{s9\_best\_007} & 203 & 40 & 8 & 16\\
\bottomrule
\end{tabular}
\end{center}

The normalized source therefore improves this particular obstruction
count.  Deleting state 74 still gives a rejected 245-state experiment.
Neither the smaller count nor the 202-pair source is evidence of a valid
machine with fewer than 246 states.

\subsection{Retargeting relative controls and preserving the bootstrap}

The next extension recomputes a break or continue operand after a rotation.
For each occurrence, the search enumerates available operands and resolves
the resulting destination in the candidate's own decoder.  It selects an
operand only if the destination equals the prescribed successor
occurrence.  Increment and decrement operands are unchanged, and both
decrement guards retain their prescribed targets.  The implemented search
uses existing break depth 1 and continue depths 1--3; it does not allocate
a new handler for a larger operand.

For physical emission, candidate control rows are drawn from already
observed combinations of control kind, operand and four-bit entry suffix.
Every occurrence of a physical row must admit the chosen row.  This is
again a finite proposal mechanism, followed by the complete symbolic
runtime audit.  The operational-control theorem in
Section~\ref{sec:control-readdressing} supplies the independent source proof
when the literal operands differ.

The retained experimental control \texttt{s10\_best\_004} changes 115
continues from depth 2 to depth 3, preserves the full resolved source graph,
and passes all 8,391 binary frame contracts at $d\geq64$.  Its blank
bootstrap has not been calibrated.  In particular, the finite survey of
30,000,000 moves from its original initial state did not find a standard
empty-root frame satisfying the inherited dirty-start conditions.  This
negative finite search is not a proof of eventual failure, and its runtime
certificate is not a substitute for a bootstrap proof.  The table is kept
only under \texttt{experiments/logic/controls}.

An additional retargeting run froze all 264 rows encountered during the
known blank prefix, in addition to the ordinary borrowed-row constraints.
Its union contained 286 protected rows.  Over 79,969 proposals, 5,234 exact
source candidates survived, but none changed a control operand or reduced
the best state count.  The obstruction exposed by this run is specific to
that frozen-prefix search family.

\subsection{Reproducible search versus deterministic replay}

The delivered scripts under \texttt{code/logic} include the color search,
normalization, control retargeting, seed generation and their small source
parsing dependencies.  Inputs and historical run records are under
\texttt{experiments/logic}.  Random seeds and measured finite counts are
retained.  The original search budgets were wall-clock limits, so rerunning
a search on another computer can change the number of proposals reached
before its deadline; a seed alone is not a claim of byte-identical replay.

The selected result also has a fixed-workload replay.  Running
\path{python -B code/logic/replay_selected.py} recomputes the rebased
physical roles, verifies all source guards, and emits exactly the 492
literal rows of the principal table.  This deterministic emission check is
separate from the independent source, complete blank-frame and universal
runtime checks.  The experiments' README gives the corresponding commands
and identifies every uncalibrated control explicitly.
